% Problem 3 — Algebraic Combinatorics (Lauren Williams)
% Source: First_Proof.tex, https://github.com/1stproof/batch-1
% Shared preamble for all problems from First_Proof.tex
% Source: https://raw.githubusercontent.com/1stproof/batch-1/refs/heads/main/First_Proof.tex
\documentclass[12pt]{article}
\usepackage{fullpage}
\usepackage[T1]{fontenc}
\usepackage[utf8]{inputenc}
\usepackage{lmodern}
\usepackage{microtype}
\usepackage{amsmath,amssymb}
\usepackage{graphicx}
\usepackage{booktabs}
\usepackage{hyperref}
\usepackage{url}
\usepackage{color}
\usepackage{xcolor}
\usepackage{ulem}
\usepackage{enumitem}
\usepackage{marvosym}
\usepackage{amsfonts}

\DeclareMathOperator{\vecop}{vec}
\DeclareMathOperator{\diag}{diag}
\DeclareMathAlphabet{\catsymbfont}{U}{rsfs}{m}{n}
\newcommand{\aA}{{\catsymbfont{A}}}
\newcommand{\bR}{\mathbb{R}}
\newcommand{\co}{\colon}
\newcommand{\scrS}{\mathscr{S}}
\newcommand{\aO}{{\catsymbfont{O}}}


\title{Problem 3: Markov Chain with Interpolation Macdonald Stationary Distribution}
\author{Lauren Williams (Harvard University)}
\date{}

\begin{document}
\maketitle

Let $\lambda=(\lambda_1 > \dots > \lambda_n \geq 0)$ be a partition with distinct parts.  Assume moreover that $\lambda$ is {\it restricted}, in the sense that it has a unique part of size $0$ and no part of size $1$.

Does there exist a nontrivial Markov chain on $S_n(\lambda)$ whose stationary distribution is given by
$$\frac{F^*_{\mu}(x_1,\dots,x_n; q=1,t)}{P^*_{\lambda}(x_1,\dots,x_n; q=1,t)} \text{ for }\mu\in S_n(\lambda)$$
where $F^*_{\mu}(x_1,\dots,x_n; q,t)$ and $P^*_{\lambda}(x_1,\dots,x_n;q,t)$ are the interpolation ASEP polynomial and interpolation Macdonald polynomial, respectively?  If so, prove that the Markov chain you construct has the desired stationary distribution.  By ``nontrivial'' we mean that the transition probabilities of the Markov chain should not be described using the polynomials $F_{\mu}^*(x_1,\dots,x_n; q,t)$.

\end{document}
